% Generated draft: requires mathematical and source review.
\documentclass[11pt]{article}
\usepackage[margin=25mm]{geometry}
\usepackage{fontspec}
\setmainfont{Latin Modern Roman}
\usepackage{amsmath,amssymb,mathtools,unicode-math}
\setmathfont{Latin Modern Math}
\usepackage{enumitem,needspace,xurl,hyperref}
\setlength{\parindent}{0pt}
\setlength{\parskip}{5pt}
\setlength{\emergencystretch}{3em}
\newcommand{\N}{\mathbb{N}}
\newcommand{\Z}{\mathbb{Z}}
\newcommand{\Q}{\mathbb{Q}}
\newcommand{\R}{\mathbb{R}}
\newcommand{\C}{\mathbb{C}}
\newcommand{\F}{\mathbb{F}}
\newcommand{\A}{\mathbb{A}}
\newcommand{\PP}{\mathbb{P}}
\newcommand{\E}{\mathbb{E}}
\newcommand{\eps}{\varepsilon}
\newcommand{\dd}{\,\mathrm{d}}
\begin{document}
\section*{UnsolvedMath ID: 20003160}
\textbf{Source problem code:} catalogue entry
\subsection{Definitions and mathematical statement}
A cohomogeneity-one Riemannian manifold has an isometric compact-group action with one-dimensional orbit space. A principal orbit is totally geodesic if its second fundamental form vanishes. Problem 35 asks to classify nonnegatively curved examples with at least one such principal orbit, without fixing a group diagram or classification equivalence. The supplied result treats only the standard effective join action
\[
 SO(p+1)\times SO(q+1)\curvearrowright S^{p+q+1},\qquad p,q\geq2,
\]
with principal isotropy $SO(p)\times SO(q)$ and the two standard singular isotropy groups.

In a normal unit-speed parameter $t\in[0,L]$, every invariant metric in this family has the form
\[
 g=dt^2+f(t)^2g_{S^p}+h(t)^2g_{S^q},
\]
where the sphere metrics have curvature one, $f,h>0$ on $(0,L)$, $f(0)=0$, $h(L)=0$, $h(0)>0$, and $f(L)>0$. Full smoothness means that near $0$, $f$ has a smooth odd extension with $f'(0)=1$ and $h$ a smooth even extension; near $L$, in $\rho=L-t$, $h$ has a smooth odd extension with derivative $1$ and $f$ a smooth even extension. In particular $h'(0)=f'(L)=0$ and $h'(L)=-1$, but these first-order conditions alone are not the full endpoint conditions.

The aid's criterion is
\[
 \sec_g\geq0\quad\Longleftrightarrow\quad f''\leq0,\ h''\leq0.
\]
Indeed the sectional-curvature types are $-f''/f$, $-h''/h$, $(1-f'^2)/f^2$, $(1-h'^2)/h^2$, and $-f'h'/(fh)$. Concavity and the endpoint derivatives force $0\leq f'\leq1$ and $-1\leq h'\leq0$, giving all the required signs.

For such a metric define
\[
 a=\inf\{t:f'(t)=0\},\qquad b=\sup\{t:h'(t)=0\}.
\]
The set of principal totally geodesic orbits is nonempty exactly when $a\leq b$. In that case $0<a\leq b<L$, $f$ is constant on $[a,L]$, $h$ on $[0,b]$, and the totally geodesic orbits are precisely those with $t\in[a,b]$. Equivalently the two active profile intervals do not overlap. The metric is a product of a rotational cap with the other sphere near each singular end, joined by a product collar of length $b-a$, possibly zero. The smooth even/odd and transition conditions remain required when that length is zero.

Every point of a qualifying metric lies on a zero-curvature plane: on the regular set a mixed sphere plane has curvature $-f'h'/(fh)=0$, and the cap product descriptions supply such planes at the singular orbits as well. Thus no metric in this classified family is quasipositively curved, meaning nonnegative everywhere with all planes strictly positive at some point. This sharp family classification does not classify arbitrary cohomogeneity-one diagrams.
\subsection{Short English statement}
Classify nonnegatively curved cohomogeneity-one manifolds having a totally geodesic principal orbit. The supplied result solves one fixed sphere-join family: the two smooth concave warping profiles must stop changing before the other starts. Their common constant interval is exactly the totally geodesic region, and every point still has a flat plane.
\subsection{Sources}
\begin{itemize}
\item Local UnsolvedMath registry, permanent ID $20003160$, source code catalogue entry; canonical location \texttt{aim-topology-notes.json}, \texttt{notes[247]}.
\item American Institute of Mathematics, \emph{Manifolds with nonnegative sectional curvature}, section 4, ``Manifolds of Cohomogeneity-one and Polar Actions,'' Problem 35, \url{https://aimath.org/WWN/nnsectcurvature/nnsectcurvature.pdf}. No proposer, remarks, or literature entries are supplied.
\item Attached research attempt 1 and summary give the join-family classification without external references. They are marked partial result, proved, and low novelty confidence, with an explicit machine-generated, non-peer-reviewed research-aid caveat requiring independent verification before citation.
\end{itemize}
\end{document}
